\documentclass[11pt]{article}

\usepackage[T1]{fontenc}
\usepackage{lmodern}
\usepackage[margin=1in]{geometry}
\usepackage{booktabs}
\usepackage{listings}
\usepackage{xcolor}
\usepackage{enumitem}
\usepackage{caption}
\usepackage[hidelinks]{hyperref}

% Inline code helper (escape underscores as \_ inside arguments).
\newcommand{\ct}[1]{\texttt{\small #1}}

\definecolor{codebg}{gray}{0.96}
\definecolor{kw}{rgb}{0.10,0.25,0.55}
\definecolor{cm}{gray}{0.45}
\definecolor{st}{rgb}{0.10,0.45,0.15}

\lstdefinestyle{py}{
  language=Python,
  backgroundcolor=\color{codebg},
  basicstyle=\ttfamily\footnotesize,
  keywordstyle=\color{kw}\bfseries,
  commentstyle=\color{cm}\itshape,
  stringstyle=\color{st},
  showstringspaces=false,
  breaklines=true,
  columns=fullflexible,
  frame=single,
  framesep=4pt,
  rulecolor=\color{gray!40},
  aboveskip=6pt, belowskip=6pt,
}

\captionsetup{font=small,labelfont=bf}

\title{\vspace{-1.5em}Catching Failures Before Customers Do:\\
Real Fault Injection and Recovery for a Video Chunking Pipeline}
\author{Justin Miller\\\small LanceDB --- Geneva DevEx}
\date{\today}

\begin{document}
\maketitle
\vspace{-1.2em}

\begin{abstract}
\noindent
We describe the fault-tolerance design and testing harness for the Geneva video
\emph{chunker}: the stage that expands each source video into fixed-length clips.
The chunker guarantees that no source row is ever silently dropped---every
failure becomes a typed error row under a small, stable contract. We then
validate that guarantee against the failures customers actually hit, using
\emph{real} fault injection rather than in-process mocks: transient object-store
faults are produced as genuine HTTP/TCP events by a reverse proxy in front of a
live S3-compatible endpoint (LocalStack), and mid-run worker loss is produced by
\ct{SIGKILL}ing a real refresh process and resuming from on-disk Geneva
checkpoints. A machine-checked harness asserts the observed behaviour against the
contract and a clean baseline, exits non-zero on any divergence, and renders a
report from the measured results. We report the outcomes and discuss the design
choices that make the faults real and the checks non-vacuous.
\end{abstract}

\section{Introduction}

Media-ingestion pipelines fail in messy ways: a video file is truncated, an
object store throttles a burst of reads, a spot instance is reclaimed mid-job.
The goal of this work is narrow and practical---\emph{find these failures in the
package we hand customers before they do}, and have a credible recovery story
when they occur. We scope the effort to a single, representative stage: the
video chunker, end to end.

Three classes of failure matter for this stage:

\begin{enumerate}[itemsep=1pt,topsep=2pt]
  \item \textbf{Data faults} --- deterministic and permanent: corrupt bytes, a
  truncated container, an empty or missing blob, a dangling pointer.
  \item \textbf{Transient object-store faults} --- non-deterministic: read
  timeouts, throttling (\ct{HTTP 503 SlowDown}), connection resets, and partial
  (short) reads that data poisoning cannot reproduce.
  \item \textbf{Worker loss} --- a process dies partway through a job (OOM kill,
  node failure, reclamation), and the job must resume without recomputing what
  already finished.
\end{enumerate}

Our contribution is twofold. First, an error-reporting \emph{contract} on the
chunker's output that turns every failure into an inspectable, groupable error
row instead of a silent drop. Second, a testing harness that validates the
contract against \emph{real} faults---real wire-level object-store errors and a
real killed process---and verifies recovery against a byte-for-byte baseline. The
chunker's own code path carries no fault-injection hooks; its happy path is
unmodified.

\section{Background: the Geneva chunker}

Geneva materializes derived columns and row expansions as \emph{materialized
views} (MVs) over LanceDB tables. The chunker is a \emph{one-to-many} MV: each
source video row expands into many clip rows. Such a view is created with
\emph{placeholder} rows and populated incrementally by \ct{view.refresh()};
refresh is checkpointed per fragment, so a partially-completed refresh can be
resumed and only unfinished fragments are recomputed. This property is the
foundation of the recovery story in Section~\ref{sec:recovery}.

The source \ct{videos} table is \emph{reference-only}: it stores metadata plus a
lightweight pointer (a row id), not the video bytes. On each worker the chunker
reads the blob itself directly from the source dataset:

\begin{lstlisting}[style=py]
blobs = dataset.take_blobs(blob_column, ids=[row_id])
video = blobs[0].readall()
\end{lstlisting}

\noindent
This keeps heavy byte movement cluster-side and off the client. It is also the
single I/O boundary at which transient object-store faults occur, and thus the
boundary our injection targets. In local mode the view is materialized on an
on-disk LanceDB database and refresh runs on a local Ray cluster; in enterprise
mode the same code runs against the remote Geneva runtime.

\section{The error contract}

Both chunker factories emit an output schema carrying a nullable
\ct{errors: list<string>} column. The contract has three invariants:

\begin{itemize}[itemsep=1pt,topsep=2pt]
  \item \ct{errors IS NULL} marks a fully clean row (never an empty list).
  \item \ct{clip\_bytes IS NOT NULL} marks a playable clip; consumers select
  only playable clips with this predicate.
  \item \textbf{No source row is ever silently dropped.} A whole-video failure
  yields exactly one row with null window fields and the message(s); a
  per-window failure yields that window's row with the failed value null.
\end{itemize}

Every message leads with a stable \emph{class tag}---the text before the first
\ct{:} or \ct{[}---so failures group in SQL:

\begin{lstlisting}[style=py]
SELECT split_part(errors[1], ':', 1) AS class, count(*)
FROM video_clips WHERE errors IS NOT NULL GROUP BY 1;
\end{lstlisting}

\noindent
Table~\ref{tab:classes} lists the classes. Blob reads retry transient errors with
exponential backoff, recording \emph{one} \ct{blob\_read\_failed[i/N]} message
per failed attempt, so the full retry history survives on the error row rather
than being collapsed to a single line.

\begin{table}[t]
\centering
\caption{Error classes emitted by the chunker.}
\label{tab:classes}
\small
\begin{tabular}{@{}lll@{}}
\toprule
\textbf{Class tag} & \textbf{Stage} & \textbf{Cause} \\
\midrule
\ct{pointer\_null}        & resolve & pointer column is null \\
\ct{blob\_missing}        & resolve & \ct{take\_blobs} returned nothing \\
\ct{blob\_empty}          & resolve & zero-length blob \\
\ct{blob\_read\_failed[i/N]} & resolve & read attempt $i$ of $N$ failed (transient) \\
\ct{decode\_failed}       & decode  & bytes are not a decodable video \\
\ct{no\_duration}         & decode  & no container/stream duration to window \\
\ct{skipped:max\_video\_s} & decode  & longer than the configured ceiling \\
\ct{encode\_failed}       & window  & a single window failed to encode \\
\ct{empty\_window}        & window  & zero-packet remux (no clip bytes) \\
\ct{no\_start\_frame}     & window  & clip produced but no decodable frame \\
\bottomrule
\end{tabular}
\end{table}

The retry loop is the only fault-relevant machinery in the happy path, and it is
deliberately ordinary:

\begin{lstlisting}[style=py]
for attempt in range(read_attempts):
    try:
        ds = ds_cache.get("ds") or lance.dataset(src_uri,
                                    storage_options=storage_opts)
        blobs = ds.take_blobs(blob_col, ids=[rid])
        ...
        video = bf.readall()
        break
    except Exception as e:                    # transient read: retry/log
        attempt_errors.append(f"blob_read_failed[{attempt+1}/{N}]: "
                              f"{type(e).__name__}: {e}")
        ds_cache.pop("ds", None)              # drop handle; reopen next try
        if attempt + 1 < read_attempts:
            time.sleep(read_sleep * (2 ** attempt))
\end{lstlisting}

\section{Fault injection methodology}
\label{sec:methodology}

The harness exercises the three failure classes with escalating realism. A guiding
principle throughout is that a \emph{simulated} fault (raising an exception
in-process to stand in for a 503) proves only that our own \ct{except} clause
runs; it does not prove the system survives what the object store actually does
on the wire. We therefore drive the transient and worker-loss cases against real
infrastructure.

\subsection{Data faults}

A corpus is poisoned purely through data---corrupt / truncated / empty / null
blobs, dangling and null pointers, a raw H.264 elementary stream (no container
duration), and an over-long video---and chunked through the \emph{unchanged}
pipeline. Because the faults are in the bytes, every resulting error row is
produced by the same code path a real run would take. The harness asserts each
source lands in its expected class (Table~\ref{tab:classes}).

\subsection{Transient object-store faults}
\label{sec:transient}

Transient faults are injected as real HTTP/TCP events by a small reverse proxy,
\ct{S3FaultProxy}, placed between the chunker and a live LocalStack S3 endpoint.
The chunker reads \ct{s3://} through the proxy; the corpus is written directly to
LocalStack (bypassing the proxy). The proxy forwards every request untouched
\emph{except} the ranged \ct{GET}s of Lance data files---the blob byte
reads---so dataset-open (bucket listing and manifest reads) always succeeds and
only the read under test is faulted. Four fault kinds are produced on the wire:

\begin{itemize}[itemsep=1pt,topsep=2pt]
  \item \textbf{Throttling}: a genuine \ct{HTTP 503 Slow Down} response with the
  S3 \ct{<Error><Code>SlowDown</Code></Error>} XML body.
  \item \textbf{Connection reset}: an abrupt TCP RST (\ct{SO\_LINGER 0}) mid-request.
  \item \textbf{Read timeout}: the socket stalls past the client's request timeout.
  \item \textbf{Short read}: a \ct{200 OK} with the true \ct{Content-Length} but a
  truncated body, so the client detects an incomplete read.
\end{itemize}

\begin{lstlisting}[style=py]
if kind == "throttle":       # real 503 SlowDown with S3 XML body
    self.send_response(503); self.wfile.write(SLOWDOWN_XML)
elif kind == "reset":        # real TCP RST
    self.connection.setsockopt(SOL_SOCKET, SO_LINGER, pack("ii", 1, 0))
    self.connection.close()
elif kind == "timeout":      # stall past the client request timeout
    time.sleep(stall_s); self.connection.close()
elif kind == "short_read":   # 200 + real Content-Length, truncated body
    self.wfile.write(body[: len(body) // 3])
\end{lstlisting}

Two knobs make the outcomes deterministic and the interaction legible. First, the
proxy can fault the first $N$ matching requests then pass through
(\emph{fail-$N$-then-succeed}) or fault every one (\emph{exhaust-all-retries}).
Second, the object-store client's own retry budget is capped
(\ct{client\_max\_retries=0} via \ct{storage\_options}) so that each fault
surfaces immediately and the chunker's own \ct{read\_retries} loop---the code we
are testing---is what drives recovery, rather than an opaque retry layer beneath
it. The chunker gained one small, general parameter (\ct{storage\_options},
forwarded to \ct{lance.dataset}); it contains no fault logic.

\subsection{Killed-worker recovery}
\label{sec:recovery}

Worker loss is produced by killing a real process. A real Geneva chunk refresh
runs in a child process; the parent polls the growing clips table and, once
some---but not all---videos have been materialized, sends \ct{SIGKILL} to the
child. This models an OOM kill or node failure far more faithfully than a raised
exception: a \ct{SIGKILL} cannot be caught, and it leaves whatever fragments had
already committed on disk. The parent then cleans up any orphaned Ray processes
and \emph{resumes} by reopening the same materialized view and calling
\ct{refresh()} again. Geneva loads the checkpointed fragments and recomputes only
the interrupted work. Finally the resumed table is compared, per video, against a
clean no-fault baseline.

\section{Verification}

The harness is only useful if its checks fail when reality diverges. Each case
carries an explicit, machine-checked assertion:

\begin{itemize}[itemsep=1pt,topsep=2pt]
  \item \textbf{Recover cases} (fail-$N$-then-succeed): the row must come back
  \emph{exactly} clean (\ct{observed == [clean]})---the retry genuinely worked.
  \item \textbf{Exhaust cases}: the source must yield \emph{exactly one} row, it
  must be a typed error row (\ct{blob\_read\_failed}, no playable clip), and no
  \ct{clean} row may appear---proving the failure is reported once and never
  dropped.
  \item \textbf{Recovery}: the kill must land mid-run with $0 < k < n$ videos
  checkpointed, and the resumed table must equal the baseline per video (clip
  counts \emph{and} error classes).
\end{itemize}

Any divergence appends to a failure list; a non-empty list logs each failure and
returns a non-zero exit code. To confirm the checks are not vacuous, we ran a
\emph{negative control}: expectations for one exhaust case were deliberately
inverted to expect \ct{clean}. The harness observed the real
\ct{blob\_read\_failed}, reported \ct{CHECK FAILED}, and exited non-zero---then
returned to green once the expectation was corrected.

\section{Results}

All results below are measured directly from the runs (geneva 0.14.0; LocalStack
S3 via the fault proxy; local Geneva + Ray for recovery). Overall outcome:
\textbf{PASS}.

\begin{table}[t]
\centering
\caption{Transient object-store faults (real HTTP/TCP events). ``Faults'' is the
number of proxy-injected failures; ``s'' is wall-clock seconds.}
\label{tab:transient}
\small
\begin{tabular}{@{}llrlll@{}}
\toprule
\textbf{Fault} & \textbf{Mode} & \textbf{Faults} & \textbf{Expected} & \textbf{Observed} & \textbf{s} \\
\midrule
read timeout        & recover & 1  & clean            & clean            & 2.0 \\
read timeout        & exhaust & 12 & blob\_read\_failed & blob\_read\_failed & 24.4 \\
503 SlowDown        & recover & 1  & clean            & clean            & 0.4 \\
503 SlowDown        & exhaust & 48 & blob\_read\_failed & blob\_read\_failed & 7.4 \\
connection reset    & recover & 1  & clean            & clean            & 0.1 \\
connection reset    & exhaust & 12 & blob\_read\_failed & blob\_read\_failed & 0.4 \\
partial/short read  & exhaust & 12 & blob\_read\_failed & blob\_read\_failed & 0.5 \\
\bottomrule
\end{tabular}
\end{table}

Table~\ref{tab:transient} shows every transient case matching its expectation.
The three recover cases return a clean, fully-chunked row---the chunker's backoff
absorbs the transient---while the four exhaust cases each land as a single typed
\ct{blob\_read\_failed[N/N]} row carrying the full attempt history, with no source
dropped. Data faults matched identically: \ct{good}$\to$\ct{clean} (2 clips),
\ct{garbage}$\to$\ct{decode\_failed}, \ct{too-long}$\to$\ct{skipped}.

For recovery, a 10-source corpus (16 baseline clips) was chunked in a child
process that was \ct{SIGKILL}ed after \textbf{8 of 10} videos had checkpointed.
After orphaned-Ray cleanup, a fresh process reopened the view and completed the
refresh; the resumed table matched the baseline for all ten videos---identical
clip counts and error classes (\ct{good-0..7} at 2 clips each,
\ct{garbage}$\to$\ct{decode\_failed}, \ct{too-long}$\to$\ct{skipped}). Only the
two videos left unfinished at kill time were recomputed on resume.

\section{Discussion and limitations}

\paragraph{Realism has a cost.} The transient and recovery harnesses depend on
infrastructure (a container runtime for LocalStack; a local Ray cluster) and take
tens of seconds each, so they are on-demand integration tools rather than unit
tests in CI. The unit suite instead pins the error contract with driver-callable
tests over poisoned data---fast and deterministic---while the real harnesses
cover the wire behaviour that unit tests cannot.

\paragraph{Timing.} The \ct{read timeout} exhaust case is slow (Table~\ref{tab:transient})
because each faulted attempt genuinely stalls to the client timeout; this is a
property of a faithful timeout, not overhead. The kill in the recovery run is
timing-based: the parent polls until partial progress exists before signalling,
and single-threaded refresh makes ``partial'' easy to catch deterministically.

\paragraph{Orphan cleanup.} A \ct{SIGKILL}ed child leaves its Ray workers
orphaned; the harness force-stops Ray before resuming so the resume starts a
clean cluster. This is the one piece of environment hygiene the real kill
requires.

\paragraph{From simulated to real.} An earlier iteration injected transient
faults in-process by raising synthetic exceptions at the read boundary. That
approach was removed: it proved only that our \ct{except} clause ran, and it put
fault logic in the chunker's happy path. Moving injection onto the wire kept the
chunker clean and made the 503, reset, timeout, and short-read cases genuine.

\section{Conclusion}

The video chunker turns every failure into an inspectable, typed error row and
never silently drops a source, and it resumes correctly from mid-run worker loss.
We validated both properties against real faults---genuine object-store errors on
the wire and a genuinely killed process---with machine-checked assertions that we
confirmed are non-vacuous. The result is a package we can hand customers with a
tested failure story rather than an assumed one.

\end{document}
